% ===================================================================
%  TABELLA nr. 11 — La citad e ses monuments. --- L'incendi.
%  Traduction 2026 d'apres texto/io, controlee sur texto/fr, texto/ca
%  et texto/oc. Consigne : voir texto/rm/10-tabella-01.tex.
%
%  DEUX INSTITUTIONS DE CETTE PAGE SE PARTAGENT EXACTEMENT COMME
%  DANS LA COLONNE LUXEMBOURGEOISE, ET LES DEUX LANGUES N'ONT RIEN
%  A VOIR L'UNE AVEC L'AUTRE.
%
%  Le tableau 11 luxembourgeois avait releve un couple :
%    la caisse d'epargne y est REFAITE a la maison — « Spuerkeess »,
%      un compose germanique ;
%    les pompiers y sont EMPRUNTES tels quels — « Pompjeeën », du
%      francais.
%  La meme page romanche, mille kilometres plus loin, avec un tout
%  autre grand voisin, fait le meme partage sur les memes deux
%  objets :
%    l'alinea 6 : « banca da spargn », un compose fait sur place ;
%    l'alinea 2 de l'incendie : « pumpiers », emprunte tel quel.
%
%  DEUX PETITES LANGUES, DEUX VOISINS DIFFERENTS, LE MEME RESULTAT :
%  LA CAUSE N'EST DONC DANS AUCUNE DES DEUX LANGUES. Elle est dans ce
%  que sont les deux institutions. Une caisse d'epargne, un village
%  peut la fonder lui-meme, avec ses propres gens et son propre
%  argent — le nom se compose avec les mots qu'on a. Un corps de
%  pompiers arrive du dehors deja constitue, en uniforme, avec son
%  reglement et son grade, et le nom vient avec le corps.
%
%  C'EST LE CRITERE DU TABLEAU 16 LUXEMBOURGEOIS RENDU VERIFIABLE.
%  On y disait que la ligne passe entre ce qui a un guichet et ce qui
%  n'en a pas ; on ne pouvait le verifier que dans une langue. Ici
%  deux langues sans rapport donnent la meme reponse sur le meme
%  couple, et c'est une confirmation, non une repetition.
%
%  Pour le reste, cette page est celle d'une vie que les locuteurs du
%  romanche ont menee en allemand : « borsa », « prefectura »,
%  « tribunal », « banca », « museum », « chaserna », « ospital »,
%  « seminari », « posta », « teater », « opera ». Les Grisons n'ont
%  qu'une ville, et elle ne parle pas romanche depuis le XVe
%  siecle. La seconde moitie du tableau est un incendie de theatre ;
%  c'est un incendie, en 1464, qui est reste dans la memoire de Coire
%  comme le moment ou la ville s'est rebatie en allemand. La
%  coincidence n'explique rien et ne prouve rien. Elle vaut d'etre
%  notee une fois, et une seule.
%
%  DEUXIEME INVERSION DE LA COLONNE, AU BLOC t11-c2-02-1. L'ido range
%  soldats (57), gendarmes (63), curieux (56) : la comparaison vient
%  avant l'objet du verbe. La premiere redaction romanche mettait le
%  complement d'objet d'abord — 61, 56, 57, 63 — et renvoji.py l'a
%  signalee. On a fait porter le comparatif sur le verbe plutot que
%  sur la phrase, ce qui remet les trois renvois dans l'ordre de
%  l'ido. La colonne reste a zero divergence.
% ===================================================================

%%K t11-apar-1 apar
\VUsaut{2.0mm}
\VUcentre{13.2pt}{90}{TERZA SERIA}
\VUsaut{3.0mm}
\VUfilet{16mm}
\VUsaut{5.6mm}
\VUtitre{15.0pt}{60}{TABELLA nr. 11}
\VUsaut{1.4mm}
\VUpk{11.4pt}{40}{La citad e ses monuments. --- L'incendi.}
\VUsaut{2.2mm}

%%K t11-tit-1 sub
\VUpk{10.2pt}{40}{Il suburbi.}
\VUsaut{3.0mm}

%%K t11-c1-01-1 p
1. --- Jau abitesch en ina gronda citad situada a 100 kilometers da l'Ocean, e ch'è
tuttina in \VUgras{port} \textsuperscript{(1)} da l'emprim rang. Il flum ch'la circumda ha
ina largezza da 400 meters e furma ina fitg bella rada.

%%K t11-c1-02-1 p
2. --- Tut la part da la citad ch'è sin ils \VUgras{cais} \textsuperscript{(2)} para dal tut
maritima ed industriala, e furma in \VUgras{suburbi} \textsuperscript{(3)} ch'è abità mo da
marinars e lavurants. Quels lavuran en las \VUgras{fabricas} \textsuperscript{(4)} ed en la
\VUgras{fabrica da gas} \textsuperscript{(5)}, da la quala l'aut \VUgras{chamin}
\textsuperscript{(6)} ed il \VUgras{gasometer} sa vesan da fitg lunsch.

%%K t11-c1-03-1 p
3. --- En il temp mesaun era la citad fortifitgada, ed ins chatta anc intgins rests
da ses mirs da defensiun, en spezial la \VUgras{porta} \textsuperscript{(8)} dal vegl
\VUgras{chastè fortifitgà} \textsuperscript{(7)}, flanchegià da ses dus \VUgras{turms}
\textsuperscript{(9)}, che vegnan ussa duvrads sco \VUgras{praschun}.

%%K t11-c1-04-1 p
4. --- En tut quel \VUgras{quartier}, che para fitg vegl, è mo ina construcziun
relativamain moderna: quai è la \VUgras{duana} \textsuperscript{(10)}. Ella è situada tranter
il cai e la \VUgras{vietta} \textsuperscript{(11)} che maina a la veglia \VUgras{chasa
communala} \textsuperscript{(12)}. Ins enconuscha facilmain quest ultim monument pervi dal
\VUgras{clutger} \textsuperscript{(13)} gigantic che domina la citad antica.

%%K t11-c1-05-1 p
5. --- Da l'autra vart da la via \VUgras{stat} in edifizi construì en il medem stil sco
la duana: quai è la \VUgras{Borsa} \textsuperscript{(14)}. Ina da sias fatschadas dat sin
ina pitschna plazza, nua ch'è la \VUgras{prefectura} \textsuperscript{(15)}. Propi è nossa
citad, senza esser ina metropola, tuttina ina impurtanta chapitala da departament.

%%K t11-tit-2 sub
\VUsaut{5.0mm}
\VUpk{10.2pt}{40}{La part la pli auta da la citad.}
\VUsaut{3.4mm}

%%K t11-c1-06-1 p
6. --- La guarnischun, bastantamain numerusa, è alloschada en vastas \VUgras{chasernas}
\textsuperscript{(16)} fitg sanas; ellas sa stendan fin al grigl dal \VUgras{parc
communal} \textsuperscript{(17)}. Il \VUgras{boulevard} \textsuperscript{(19)} che maina a
quest
parc passa davos la \VUgras{banca da spargn} \textsuperscript{(18)} e va fin a la plazza da la
\VUgras{catedrala} \textsuperscript{(20)}.

%%K t11-c1-07-1 p
7. --- Tut quels monuments sa stendan sco in amfiteater sin ina collina, nua ch'è era
noss vast \VUgras{gimnasi} \textsuperscript{(21)}.

%%K t11-c1-07-2 p
Sch'ins va giu tras ina via in pau pendenta, che maina a la Via gronda, arriva ins
al \VUgras{Tribunal} \textsuperscript{(22)}, avant il qual sa stenda in agreabel \VUgras{iert
public} \textsuperscript{(26)}. Quella \VUgras{plazza da l'iert} n'è betg fitg gronda, ma
ella fa en il mez da la citad ina oasa da verd e da frestgezza. Perquai è ella fitg
frequentada dals citadins, ch'han gugent da vegnir a seser sut la tenda dal
\VUgras{chaffè} \textsuperscript{(25)}, per ascultar ils musicists schuldads che sunan la
gievgia e la dumengia en il \VUgras{kiosc} \textsuperscript{(27)} mess tranter ils praus ed
ils massivs.

%%K t11-c1-08-1 p
8. --- Sin in da quels praus stat ina \VUgras{statua} \textsuperscript{(28)} da bronz,
monument erigì per la memoria d'in da noss concitadins, ch'ha fatg gronds servetschs
a sia citad da naschientscha.

%%K t11-tit-3 sub
\VUsaut{4.0mm}
\VUpk{10.2pt}{40}{Ils meds da transport.}
\VUsaut{3.4mm}

%%K t11-c1-09-1 p
9. --- Enfin ussa n'è nossa citad anc betg illuminada cun la glisch electrica, ella
ha mo \VUgras{brischaders da gas} \textsuperscript{(29)}. Ma las \VUgras{gasettas localas} fan
campagna per ch'quel sistem d'illuminaziun vegnia meglierà, \VUgras{sco} ch'ins ha gia
perfecziunà \VUgras{il sistem da tracziun dals} trams. \VUgras{Ils} vegls chars,
\VUgras{tratgs} da \VUgras{chavals}, èn praticamain vegnids remplazzads cun \VUgras{trams
electrics} \textsuperscript{(31)}, che transportan svelt ils viagiaders d'ina fin da la
\VUgras{citad a l'autra}, per in pretsch fitg \VUgras{bunmartgà}.

%%K t11-c1-10-1 p
10. --- Sper il Tribunal è la \VUgras{Banca} \textsuperscript{(32)}, nua ch'ins escumpta
las valurs commerzialas. Ils \VUgras{incassaders da la banca} \textsuperscript{(33)} van a
retscheiver a chasa ils debits.

%%K t11-tit-4 sub
\VUsaut{4.0mm}
\VUpk{10.2pt}{40}{Art ed instrucziun.}
\VUsaut{3.4mm}

%%K t11-c1-11-1 p
11. --- Il bel edifizi ch'è dasperas è il \VUgras{museum}. El cuntegna ensemen la
\VUgras{biblioteca} \textsuperscript{(34)} da la citad, las bellas \VUgras{collecziuns da
scienzas natiralas} \textsuperscript{(35)} (museum da natira) ed in grazius \VUgras{museum da
picturas} \textsuperscript{(36)}, ch'è ina splendida \VUgras{exposiziun} permanenta.

%%K t11-c1-11-2 p
El vegn avert per il public duas giadas l'emna, ma ils esters dastgan al visitar
mintga di per ina pitschna dunativa al \VUgras{guardian} \textsuperscript{(37)}. Ils
\VUgras{picturs} \textsuperscript{(38)} vegnan là per lavurar. Ins als vesa \VUgras{avant}
lur cavalet, cun ina \VUgras{paletta} en maun, copiond las picturas famusas.

%%K t11-c1-12-1 p
12. --- Sin l'autezza, davos il museum, vesa ins la \VUgras{cupola} \textsuperscript{(40)}
da l'\VUgras{ospital} \textsuperscript{(39)} e las construcziuns dal \VUgras{seminari}
\textsuperscript{(41)}, en il qual vegnivan scolads ils magisters da nossas scolas
communalas.

Sin la plazza dal Teater è ina \VUgras{posta} \textsuperscript{(43)} installada en ina
\VUgras{chasa privata} \textsuperscript{(42)}.

%%K t11-tit-5 sub
\VUsaut{5.0mm}
\VUpk{11.4pt}{40}{L'incendi.}
\VUsaut{3.0mm}

%%K t11-tit-6 sub
\VUpk{10.2pt}{40}{Clom d'incendi.}
\VUsaut{3.4mm}

%%K t11-c2-01-1 p
1. --- Ina nova terrorisanta sa derasa tras la citad: in incendi è gist stà en il
teater! En il mument ch'jau arriv sin la \VUgras{plazza} \textsuperscript{(96)}, è la fulla
gia radunada sin il \VUgras{trottuar} \textsuperscript{(46)} e sin la \VUgras{chaussée}
\textsuperscript{(47)}. Ils \VUgras{impiegads da la posta} \textsuperscript{(44)} ed ils
\VUgras{purtaders da brevs} \textsuperscript{(45)} sortan da la posta. In \VUgras{manader dal
tram} \textsuperscript{(48)} ha stuì fermar ses char per laschar passar l'\VUgras{ambulanza}
\textsuperscript{(50)} da la citad, ch'arriva cun in \VUgras{chirurg}
\textsuperscript{(52)} ed in \VUgras{sanitari} \textsuperscript{(51)} per tgirar in
\VUgras{blessà} \textsuperscript{(53)} purtà sin ina \VUgras{brancarda}
\textsuperscript{(54)} sper il \VUgras{kiosc da gasettas} \textsuperscript{(55)}.

%%K t11-c2-02-1 p
2. --- Ina cumpagnia d'infantaria, ch'era sin via da return da l'exercizi, s'ha
fermada per assicurar il servetsch da l'urden cun ils \VUgras{schandarms a chaval da la
citad} \textsuperscript{(61)}. \VUgras{Ils chavalgiaders}, dals quals ils chavals sa
sblatgan, san meglier ch'ils \VUgras{schuldads} \textsuperscript{(57)} u ch'ils
\VUgras{schandarms} \textsuperscript{(63)} far turnar enavos ils \VUgras{curius}
\textsuperscript{(56)}.
Dal rest èn ils schandarms fitg occupads. In d'els mussa als \VUgras{polizists}
\textsuperscript{(64)} nua ch'ins sto transportar in auter \VUgras{accidentà}
\textsuperscript{(65)}, in curaschus pumpier crudà giu d'ina stgala.

%%K t11-c2-03-1 p
3. --- L'\VUgras{uffizier} \textsuperscript{(66)} precisescha il lieu nua ch'ins sto
metter la \VUgras{pumpa a vapur} \textsuperscript{(67)} ch'arriva. Spetgond quella pussanta
maschina, ha el fatg installar en pressa ina \VUgras{pumpa a maun} \textsuperscript{(68)},
da la quala il lung \VUgras{tub} \textsuperscript{(69)} bitta aua sco in torrent sin il
fieu.

Sis pumpiers la manovreschan, entant che lur cumpogn, che tegna la \VUgras{lantscha}
\textsuperscript{(70)}, dirigia il \VUgras{sgrizzel} \textsuperscript{(71)} vers il center da
l'incendi.

%%K t11-c2-04-1 p
4. --- Da l'autra vart dal teater han ins pusà la gronda \VUgras{stgala}
\textsuperscript{(72)}. Ella pussibilitescha als pumpiers da s'avischinar a las flommas
che sprizzan tranter vortischs da nair \VUgras{fim} \textsuperscript{(75)}.

%%K t11-tit-7 sub
\VUsaut{5.0mm}
\VUpk{10.2pt}{40}{Bravas acziuns.}
\VUsaut{3.4mm}

%%K t11-c2-05-1 p
5. --- L'\VUgras{incendi} \textsuperscript{(74)} è naschì en il foyer dal \VUgras{teater}
\textsuperscript{(73)}, entant ch'ins repetiva l'\VUgras{opera}, da la quala l'emprima
\VUgras{represchentaziun} avess gì lieu damaun. Per fortuna eran mo pauc
\VUgras{spectaturs} \textsuperscript{(76)} en la sala. Ils pauprs sa èn refugiads sin il
\VUgras{balcun} \textsuperscript{(77)}, nua che curaschus \VUgras{pumpiers}
\textsuperscript{(86)} als vegnan a socurrer cun stgalas e cordas.

%%K t11-c2-06-1 p
6. --- Sut il \VUgras{peristil} \textsuperscript{(78)}, nua ch'l'\VUgras{affisch}
\textsuperscript{(79)} infurmescha davart la represchentaziun, vesa ins ils
\VUgras{musicists} \textsuperscript{(81)} da l'\VUgras{orchester} \textsuperscript{(80)} fugir,
purtond cun els lur instruments: \VUgras{violina} \textsuperscript{(81)}, \VUgras{flauta}
\textsuperscript{(82)}, \VUgras{violoncell} \textsuperscript{(83)}, \VUgras{trumbun}
\textsuperscript{(84)}, \VUgras{tambur} \textsuperscript{(85)}, etc. Ins sa dat fadia da
salvar ina part da la \VUgras{mobiglia} \textsuperscript{(87)}.

%%K t11-c2-07-1 p
7. --- Ils \VUgras{acturs} \textsuperscript{(89)} e las \VUgras{actricas}
\textsuperscript{(88)}, ils \VUgras{chantadurs} e las \VUgras{chantadras} han stuì fugir cun
lur \VUgras{custum} \textsuperscript{(90)} da teater. Il tenor explitgescha ad in
\VUgras{schurnalist} \textsuperscript{(92)} co che quai è capità. \VUgras{Il reporter} è
currì immediatamain dal emprim mument cun las autoritads: \VUgras{prefect}
\textsuperscript{(91)}, \VUgras{sindic} \textsuperscript{(93)}, \VUgras{commissari da polizia}
\textsuperscript{(94)} ed \VUgras{uffizial da la giustia} \textsuperscript{(95)}, che vegnan
ad inquirir per giuditgar davart las responsabladads.

\VUsaut{6.0mm}
\VUfilet{22mm}
